\section{Results}
\label{sec:results}

Table~\ref{tab:battery} lists every measurement reported below together with the control it is read
against. The order follows the question rather than the chronology: first whether the data
determine the internal quantity at all, then what the network does with the freedom that remains,
then what that costs the thermodynamics.

\subsection{The observable determines two of fifty directions of the profile}
\label{sec:ident}

A $\sigma$-profile is a point in a $51$-dimensional space, and normalisation fixes one direction, so
a profile can move in $50$. The question is how many of those the closure exposes. For each
solvent we collected the rows in which it appears, differentiated the predicted activity
coefficient with respect to its profile, and took the spectrum of the resulting information matrix.
Over $25$ solvents and five seeds, the number of directions carrying $99\%$ of the information is
$2$, with a range of $1$ to $5$; the median condition number is $1.4\times10^{33}$. The remaining
directions are not merely poorly determined. They are flat to numerical precision, so a profile can
move along them without changing any prediction the closure makes.

The same calculation on random profiles of equal area returns the same count. The deficiency is
therefore structural, a property of the exchange kernel and of the diversity of partners available,
rather than something the trained head produced. That is the stronger of the two readings: it holds
wherever the profile is linearised, not only at the point this network happens to occupy.

The crystal term is underdetermined in the same way and for a reason that is easier to picture. The
enthalpy of fusion and the temperature slope of the activity coefficient both enter
equation~\eqref{eq:sle} through $1/T$, so a fit can trade one against the other. For the
four-parameter fit the condition number of the information matrix is $2.3\times10^{5}$, and no
estimator can determine the enthalpy of fusion from a solubility series to better than about
$\pm20$~kJ\,mol$^{-1}$, which is wider than the range the quantity takes across these compounds.
Supplying the melting temperature and the enthalpy of fusion as independent labels reduces that
bound to $\pm1.7$~kJ\,mol$^{-1}$.

What the network places in the undetermined directions is not decided by the data. The rest of this
section measures what it places there.

\subsection{The network places charge density where the reference carries none}
\label{sec:windows}

The reference tabulation is sparse. A molecule with no hydrogen-bond donor is tabulated with no
cavity surface carrying donor-side charge density, and $58\%$ of all reference cells across the
compounds examined here are exactly zero, a median of $29$ bins of $51$ per molecule.

The learned profiles are not sparse anywhere. Figure~\ref{fig:windows} shows the donor-window area
for sixteen common solvents. Twelve of the sixteen are exactly empty in that window in the
reference; all sixteen carry area there in the learned profile, between $0.1\%$ and $5.2\%$ of the
molecule's surface. The hydrogen-bonding term reads that area regardless of how small it is,
because a segment qualifies for the term by its charge density alone.

Where the misfit lands is not uniform across the charge-density axis. Take the residual between the
learned and the reference profile over the $71$ molecules that appear in a labelled test row, and
express it relative to the share of surface area each window holds. The donor window is then
over-represented by a factor of $4.0$ and the acceptor window by $2.0$, while the nonpolar core is
under-represented at $0.88$. The profiles a molecule receives as solvent give $3.8$, $2.0$ and
$0.89$ for the same three windows, so the orientation of the displacement does not depend on which
role the molecule plays. The network's departure from the tabulation concentrates in the regions
that switch on the model's hydrogen-bonding term and thins in the region that does not.

\subsection{The output layer is not the cause}
\label{sec:pedestal}

An obvious explanation is architectural. The head ends in a normalisation that cannot return an
exact zero, so a profile that should be empty in a bin will instead be slightly populated there,
and the defect would be an artefact of the parameterisation. The explanation is testable, and it
fails.

Figure~\ref{fig:pedestal}a shows, for every bin in which the reference is exactly zero, how much
mass the trained head places there. The distribution reaches $1.0\times10^{-11}$, and between
$34\%$ and $51\%$ of those bins already lie below $10^{-5}$. A layer that can reach $10^{-11}$ is
not what holds the remaining bins at $10^{-5}$. Figure~\ref{fig:pedestal}b makes the same point
from the other direction. An untrained network of the same architecture places a median of $0.571$
of the profile's mass in the bins where the reference is empty; after training the figure is
$0.022\pm0.004$ over five seeds and both roles. The residual population is something training
arrives at, not something the parameterisation imposes.

The two roles agree. The pedestal is $0.024\pm0.004$ on the profiles a molecule receives as solute
and $0.020\pm0.003$ on those it receives as solvent, a difference of $0.004$ against a
seed-to-seed spread of the same size. This measurement and the one in
Section~\ref{sec:windows} therefore hold separately for each of the two profiles a molecule is
given, and do not depend on the fact that it is given two.

The point extends beyond the present pipeline. A recent physics-embedded model uses a rectified
output with the same normalisation and can therefore emit exact zeros \cite{yang2025tennet}. The
measurement above indicates that the difference between the two output layers is not what decides
whether a learned profile resembles the tabulated one.

\subsection{A molecule is not ideal in itself}
\label{sec:refstate}

Any activity coefficient model must return unity for a pure component, since a molecule surrounded
by copies of itself is in its reference state by definition. In COSMO-SAC the property holds
identically: the residual term is a difference between a segment's activity in the mixture and in
the pure compound, and the two coincide when one profile is used on both sides. Supplying a single
profile to both sides returns exactly zero here, which confirms that the implementation carries the
property.

The pipeline does not. Because the head reads a role-dependent representation, a compound receives
one profile as solute and a different one as solvent, so pairing a molecule with itself pairs two
different profiles. The median departure from ideality is between $0.012$ and $0.044$ in
$\ln\gamma$ across the five seeds, the ninety-fifth percentile between $0.09$ and $0.34$, and the
largest single value is $0.80$. The violation belongs to the learned intermediate rather than to
the thermodynamic relation. It is also invisible in any evaluation that never pairs a compound with
itself, which is every evaluation on solubility or on mixture data.

\subsection{What misplaced density costs, and where}
\label{sec:cost}

The preceding measurements establish that density sits where the reference says it should not. They
do not say what that costs, and the cost is not uniform. Figure~\ref{fig:costmap}a reports, for
each bin, the mean change in the predicted activity coefficient produced by moving a fixed amount
of mass from the core into that bin. Moving mass within the core is almost free. Moving the same
mass into the outer windows costs $21$ times as much as a median over seeds, with a range of $21$
to $30$; computed over the rows used for scoring rather than for training the ratio is $15$, with a
range of $13$ to $25$. Random profiles of the same total area give $1.6$ and $1.5$, so the
structure reflects the chemistry of the exchange rather than the geometry of the grid.

The cost does not grow smoothly with distance along the axis. It is flat across the core, rises
steeply as the hydrogen-bonding threshold is crossed, and then saturates. Once mass is inside the
bonding region, pushing it further out changes the prediction comparatively little. The consequence
is shown in Figure~\ref{fig:costmap}b. Reading the measured cost as a potential and taking its
gradient gives the weight under which the usual transport distance between two profiles would price
their difference correctly, and that weight peaks at the boundary rather than deep in the outer
window. A treatment that weighted the far corner most heavily, which is what the geometry of the
grid would suggest, would weight the wrong place.

\subsection{The intermediate does not decompose}
\label{sec:nonadditive}

If the learned profile were a $\sigma$-profile carrying a per-bin error, then correcting part of it
should help in proportion to the part corrected. It does not. Replacing the learned shape with the
reference shape over the whole grid improves the predicted activity coefficients on all five seeds.
Replacing it over a randomly chosen third of the bins makes the prediction worse than leaving the
profile entirely alone in $20\%$ of draws, taken over twenty random thirds at each of five seeds and
ranging from $5\%$ to $30\%$ by seed. Recovery of the gap across those draws runs from $-213\%$ to
$+152\%$.

A profile assembled from correct values in some bins and learned values in the rest is not
intermediate between the two. The segment-activity calculation is a self-consistent solution over
the whole distribution, so an internally inconsistent profile can sit far from both of the profiles
it was built from. Two things follow. The first is methodological: the share of the error
attributable to any subset of bins is not defined on this surface, so an analysis that partitions
the grid and attributes error to the parts cannot be carried out, whatever the partition. The
second concerns the premise under test. Partial correctness of the intermediate is not partially
useful, and is frequently worse than consistent incorrectness.

\subsection{The error that reaches the thermodynamics is in the occupied bins}
\label{sec:attribution}

One partition is well posed, because both of its parts are substituted in full: the bins where the
reference is nonzero, and the bins where it is zero. Figure~\ref{fig:attribution}a reports both.
Supplying the reference values in the bins the reference occupies recovers $89\%$ of the gap
between the learned and the reference profile, positive on all five seeds and ranging from $68\%$
to $109\%$. Setting exact zeros in the bins the reference leaves empty recovers $-9\%$ on average,
with a per-seed range from $-36\%$ to $+34\%$, and the result sits inside the spread of a control
that zeroes the same number of randomly chosen cells.

The spurious density is therefore not what the thermodynamics responds to. The error that
propagates through the closure lives in the magnitude of the profile where the profile is
populated. The cost map of Section~\ref{sec:cost} does not localise this error either.
Figure~\ref{fig:attribution}b shows that substituting the reference into the third of the grid the
closure is most sensitive to recovers no more of the gap than substituting into a randomly chosen
third, and that the sign varies across seeds. Sensitivity per unit of displaced mass and total
contribution to the error are different quantities, and the core carries most of the mass.

The profile's shape and the molecule's cavity area, which COSMO-SAC consumes as a separate scalar,
split the gap unevenly. Replacing the whole shape while keeping the learned area moves the mean
error from $1.75$ to $1.08$, against $0.76$ when the reference profile is supplied with its own
area, so the shape accounts for $0.68$ of the gap and the area for the rest. This is a
decomposition of two full substitutions rather than an attribution over bins, and is not subject to
the objection raised in Section~\ref{sec:nonadditive}.

\subsection{Supplying the reference improves one quantity and degrades the other}
\label{sec:twosided}

The substitution the premise licenses can be scored on two quantities at once: the activity
coefficient the profile exists to produce, and the solubility the pipeline is built to predict.
Both are measured on the same checkpoints.

On the activity coefficients, replacing the learned profiles with the reference profiles moves the
mean absolute deviation from $1.751\pm0.063$ to $0.765$, an improvement on every one of the five
seeds and more than fourteen times the seed-to-seed spread of the arm. On the solubility, the same
replacement at prediction time moves the error in the opposite direction, by
$0.431\pm0.044$ in mean absolute error over the $5274$ test rows in which a reference profile was
available for the solvent, with per-seed values of $0.456$, $0.381$ and $0.457$.

The control is internal to the same comparison. No reference profile exists for either component in $37$ test
rows, so nothing is substituted in them, and across those rows the two arms differ by at most
$4\times10^{-6}$ in mean absolute error. The effect is therefore carried by the rows in
which the substitution actually happened, and not by any difference in how the two arms were
scored. The substitution is a solvent-side substitution throughout: of the $147$ distinct solutes
in the test rows only seven are tabulated, so no row exists in which the solute profile alone is
replaced. What is reported is what the available reference data permit, and the solvent channel is
where the comparison lives.

The two directions are the result this work was undertaken to obtain. A better value for the
internal quantity makes the thermodynamic model it feeds more accurate and makes the pipeline less
accurate.
